\section{Multi-modal Multi-task Masked Autoencoders(MultiMAE)}
Multi-modal Multi-task Masked Autoencoders (MultiMAE) \cite{multimae2022} extends
the MAE framework (RGB-only) to handle diverse modalities such as RGB image,
depth maps and semantic segments simultaneously.

\subsection{Multimodal Input Handling}
In MultiMAE, each modality uses its own specific linear projection to project raw
inputs into a set of tokens that all share the same unified latent dimension.
The length of th token sequence can differ between modalites depending on therir
original structure and how they are tokenized. Once projected, these variable-length
sequences are concatendated into a single unified sequence for the encoder.

\subsection{Cross-Modal Predictive Coding}
To help the model distinguish spatial locations and modalities, 2-D positional embeddings
and modality-specific learnable embeddings are added to the token representations.
MultiMAE masks patches across both spatial locations and input modalities. This forces
the network to learn deep inter-modal correlations, allowing it to "hallucinate"
or reconstruct missing sensors from whatever subset of data is available.

\subsection{Symmetric Dirichlet Sampling}
MultiMAE uses a Dirichlet distribution to randomly sample the number of tokens
contributed by each modality during pre-training. It determines the proportions of
visible tokens $(\lambda_{RGB}, \lambda_{D}, \lambda_{S})$ from each modality,
ensuring that the model learns to handle any combination of available inputs.

\begin{equation}
    (\lambda_{RGB}, \lambda_{D}, \lambda_{S}) \sim \text{Dir}(\alpha)
\end{equation}

Subject to the unity constraint:
\begin{equation}
    \sum_{i \in \{RGB, D, S\}}\lambda_{i}= 1, \quad \lambda_{i}\geq 0
\end{equation}
Where the concentration parameter $\alpha > 0$ determines the sampling diversity:
\begin{itemize}
    \item $\alpha = 1$: Equivalent to a uniform distribution over the simplex, which
        is default and exposed to all possible "sensor failure" scenarios during
        pre-training to force model to learn how ot reconstrcut any modaility from
        whatever sparse subset of data is available.

    \item $\alpha \ll 1$: Modality-specific "extreme" sampling where most tokens
        are drawn from a single stream.

    \item $\alpha \gg 1$: Leads to an equal distribution of tokens across all
        modalities, which ensuring a fixed, identical number of tokens from each
        modality.
\end{itemize}




\section{BP4D+ Multimodal Spontaneous Expression Database}
The experimental data for this research is sourced from the BP4D+ Multimodal Spontaneous
Emotion Corpus \cite{zhang2013high}, a comprehensive dataset featuring
synchronized sequences from 140 subjects. To ensure the collection of authentic affective
states, the dataset utilizes an emotion elicitation protocol consisting of ten
distinct tasks, such as humor, physical threat, and cold pressor tests.

The corpus provides a diverse range of high-resolution modalities.
\begin{itemize}
    \item \textbf{Visual Data:} 2D texture images, 3D facial models, and thermal
        video sequences, all recorded at a frame rate of 25 fps.

    \item \textbf{Physiological Signals:} Continuous 1D time-series
        data—including blood pressure, heart rate, respiration rate, and electrodermal
        activity (EDA)—captured at a sampling rate of 1000 Hz.

    \item \textbf{Derivatives and Annotations:} Expert-coded Facial Action
        Coding System (FACS) intensity and occurrence, as well as tracked head poses
        and feature points for the 2D, 3D, and thermal streams.
\end{itemize}

This multi-sensory synchronization makes the BP4D+ an ideal foundation for
developing models capable of learning deep cross-modal correlations.

\section{Challenges of Video-based Multimodal Data}
Since the BP4D+ dataset is a video-based multimodal dataset, it presents several
challenges for masked autoencoding that go beyond the static image setting
discussed in the previous sections. This section highlights the difficulties most
relevant to this work.

\subsection{Temporal Synchronization of Heterogeneous Data}
The BP4D+ corpus records its modalities at fundamentally different sampling
rates: the 1D physiological signals are captured at $1000\,\mathrm{Hz}$, whereas
the 2D, 3D, and thermal visual streams are recorded at only $25\,\mathrm{fps}$. This
temporal mismatch must be resolved before the data can be processed jointly.

A first challenge is therefore to align the 1D physiological time-series with the
2D and thermal video frames in a synchronized manner, ensuring that
corresponding time steps refer to the same physical moment in the recording.

\subsection{Explosive Token Count in Joint Space-Time Cube Embeddings}
In the Joint Space-Time cube embeddings there will be temporal dimension added. As
example, a 16-frame clip represented as raw input tensor
$(Batch Size, Channel s, 16, Height, Width)$. After it got tokenized into $2\times
16 \times 16$ cubes and mapped into the latent dimension D (768 for ViT-B), its
shape become $(Batch Size, D, 8, H/16, W/16)$. The token length of a 16-frame clip
is 8 times as in a single image $(H/16 \times W/16)$. Because the self-attention
mechanism scales quadratically $O(N^{2})$ with the number of tokens, it leads to
a $64 \times$ increase in the computation load.

\subsection{Information Leakage and Temporal Tube Masking}
A critical challenge in video signals is information leakage. During pre-training,
a model can easily ``cheat'' during reconstruction by copying pixels from unmasked
adjacent frames rather than learning high-level spatiotemporal structures. This
is particularly problematic for video, where consecutive frames are highly correlated
and provide a trivial shortcut to low reconstruction error.

Tube masking mitigates this issue by enforcing a consistent spatial mask across the
entire temporal axis. By masking entire ``tubes'' through time, the model is prevented
from copying spatial regions from neighboring frames. This forces the attention
mechanism to reason over high-level spatiotemporal structures instead of relying
on low-level temporal correspondences.
